% !TeX root = ../../tesis.tex
%%%%%%%%%%%%%%%%%%%%%%%%%%%%%%%%%%%%%%%%%%%%%%%%%%%%%%%%%%%%%%%%%%%
% SECCIÓN 5.3.4 — FACTOR DE DESVIACIÓN (DI)
%%%%%%%%%%%%%%%%%%%%%%%%%%%%%%%%%%%%%%%%%%%%%%%%%%%%%%%%%%%%%%%%%%%
\subsection{Factor de Desviación (\texorpdfstring{$DI$}{DI})}
\label{subsec:factor-desviacion}

Cuando alguien planea un viaje con el mapa en mano, la distancia en línea recta
entre origen y destino parece razonable. Lo que la red de transporte impone en
la práctica es otra cosa: rutas que obligan a converger hacia el centro antes de
salir hacia la periferia opuesta, trayectos que multiplican el recorrido
respecto al camino ideal. El \textbf{Factor de Desviación} ($DI$) cuantifica
exactamente esa brecha: es la razón entre la distancia real recorrida sobre la
infraestructura —incluyendo la primera y última milla peatonal— y la distancia
teórica en línea recta entre el mismo par origen-destino.

\begin{equation}
    DI = \frac{D_{\text{red}} + D_{\text{caminata}}}
              {D_{\text{haversine}}(O, D)}
    \label{eq:factor-desviacion}
\end{equation}

\noindent donde:
\begin{itemize}
    \item $D_{\text{red}}$: distancia acumulada sobre el trazo real de
          la ruta óptima en tiempo.
    \item $D_{\text{caminata}}$: suma de la primera y última milla
          peatonal (acceso y egreso a la red).
    \item $D_{\text{haversine}}(O, D)$: distancia en línea recta entre
          origen $O$ y destino $D$, calculada con la fórmula de
          Haversine.
\end{itemize}

Un valor de $DI = 1.0$ indica una ruta perfectamente directa; valores superiores
a $1.5$ revelan ineficiencias topológicas significativas en el par
origen-destino analizado.

% ---------------------------------------------------------------
\subsubsection{Distribución en la Red}
\label{subsubsec:di-distribucion}
% ---------------------------------------------------------------

Sobre una muestra de \textbf{100 pares origen-destino} (semilla reproducible
$= 42$), el modelo arrojó los estadísticos globales del
Cuadro~\ref{tab:di-estadisticas}. La distribución resultante —con media de $1.701$ y
valores extremos entre $1.03$ y $4.42$— confirma que el recorrido promedio en la
red metropolitana supera en cerca del 70\,\% la distancia ideal.

\begin{table}[H]
    \centering
    \caption[Estadísticos globales del Factor de Desviación]{%
    Estadísticos descriptivos del $DI$ sobre 100 pares
    origen-destino de la red de la \zmvm{}.
    Corrida: 2026-10-04.}
    \label{tab:di-estadisticas}
    \begin{tabular}{|l|r|}
        \hline
        \textbf{Estadístico} & \textbf{Valor} \\
        \hline
        $n$ (rutas calculadas)     & 100   \\
        Media ($\bar{DI}$)         & 1.701 \\
        Desviación estándar        & 0.566 \\
        Mínimo                     & 1.03  \\
        Percentil 25               & 1.34  \\
        Mediana                    & 1.54  \\
        Percentil 75               & 1.90  \\
        Máximo                     & 4.42  \\
        \hline
    \end{tabular}
    \fuentefigura{Elaboración propia con \texttt{VFTModel}
    \citep{cabrera2026VFTModel}. Corrida: 2026-10-04.}
\end{table}

Estos valores incorporan la corrección del cálculo de distancias del
\texttt{VFTModel}. En versiones anteriores, la distancia acumulada a lo largo
del trazo se reiniciaba en cada vértice que caía dentro del radio de captura de
una estación; parte de las aristas medía entonces menos que la línea recta entre
sus extremos y la muestra llegaba a registrar valores de $DI$ menores que 1,
geométricamente imposibles. Ahora la distancia se acumula sobre todo el trazo y
cada estación se proyecta sobre él, de modo que ninguna arista mide menos que la
distancia en línea recta y el mínimo de la muestra es $1.03$. El Espectro MJG de
§\ref{subsec:espectro-mjg} usa una muestra mayor, de 500 pares con la misma
semilla, cuya mediana es $1.52$.

El registro completo de los 100 pares analizados —con sus distancias
individuales, valores de $DI$ y número de saltos— se incluye en el
Cuadro~\ref{tab:anx-di-corrida-completa} del
Anexo~\ref{anx:visualizaciones-analiticas}.

La Figura~\ref{fig:mapa-di} sitúa esa distribución sobre el territorio,
permitiendo leer la magnitud del rodeo como fenómeno espacial antes de
desagregarla por alcaldía.

\begin{figure}[H]
    \centering
    \includegraphics[width=0.90\textwidth]{%
        Figures/Mapas/Qgis_Resources/RedActual/Mapa_DF}
    \caption[Distribución espacial del Factor de Desviación ($DI$)
    por demarcación]{%
    Factor de Desviación ($DI$) promedio por demarcación de la \zmvm{},
    calculado sobre la muestra de 100 pares origen-destino y asignado según
    la demarcación de origen; entre paréntesis, el número de rutas de cada
    una. El gradiente verde-a-rojo señala las zonas donde la red impone
    mayores rodeos respecto a la distancia lineal. Las demarcaciones con una
    o dos rutas tienen valores solo indicativos. Corrida: 2026-10-04.}
    \label{fig:mapa-di}
    \fuentefigura{Elaboración propia con \textbf{QGIS};
    datos: \texttt{VFTModel} \citep{cabrera2026VFTModel},
    \gtfs{} \citep{adip_gtfs_2024} e INEGI 2023.}
\end{figure}

La Figura~\ref{fig:grafica-di-alcaldia} muestra la misma muestra par por par:
compara la distancia recorrida en la red con la distancia en línea recta, de
modo que la separación de cada punto respecto de la diagonal expresa la magnitud
del rodeo.

\begin{figure}[H]
    \centering
    \includegraphics[width=0.92\textwidth]{%
        Figures/Mapas/Tableau_Resources/RedActual/VFTModel/UniGrafica/%
        Grafica_VFTModel_2026_Factor_desviacion}
    \caption[Factor de Desviación --- dispersión de los 100 pares
    origen-destino]{%
    Distancia en la red frente a distancia en línea recta para los 100 pares
    origen-destino de la muestra, clasificados en cuatro categorías de
    desviación. La diagonal corresponde a $DI = 1$; cuanto más se aleja un
    punto de ella, mayor es el rodeo que impone la red.
    Corrida: 2026-10-04.}
    \label{fig:grafica-di-alcaldia}
    \fuentefigura{Elaboración propia con \textbf{Tableau};
    datos: \texttt{VFTModel} \citep{cabrera2026VFTModel}.}
\end{figure}

El Cuadro~\ref{tab:di-alcaldia} resume esa lectura por demarcación. Cada par se
asigna a la demarcación de su origen, de modo que la muestra de 100 pares se
reparte entre 20 demarcaciones y varias quedan representadas por una o dos
rutas; en esos casos el promedio es solo indicativo.

\begin{table}[H]
    \centering
    \caption[Factor de Desviación promedio por demarcación]{%
    Demarcaciones con mayor $DI$ promedio en la muestra de 100 pares
    origen-destino, asignados según la demarcación de origen. $n$: número de
    rutas. Corrida: 2026-10-04.}
    \label{tab:di-alcaldia}
    \small
    \begin{tabular}{|l|r|r|r|}
        \hline
        \textbf{Demarcación} & \textbf{$n$} & \textbf{$DI$ medio} &
        \textbf{DE} \\
        \hline
        Chimalhuacán   & 1 & 3.73 & ---  \\
        Tláhuac        & 2 & 3.26 & 1.65 \\
        Milpa Alta     & 2 & 3.13 & 0.80 \\
        Naucalpan      & 1 & 2.63 & ---  \\
        Nezahualcóyotl & 1 & 1.98 & ---  \\
        Coyoacán       & 6 & 1.84 & 0.51 \\
        Iztapalapa     & 9 & 1.80 & 0.46 \\
        \hline
    \end{tabular}
    \vspace{0.4em}

    \noindent Las 13 demarcaciones restantes registran un $DI$ medio entre
    1.37 y 1.68. Sin desviación estándar cuando $n = 1$.

    \fuentefigura{Elaboración propia con \texttt{VFTModel}
    \citep{cabrera2026VFTModel}.}
\end{table}

El patrón es consistente con los indicadores anteriores: las demarcaciones que
registran menor cobertura y menor densidad de nodos de alta Fuerza Capilar son
también las que imponen los mayores rodeos a sus usuarios.

% ---------------------------------------------------------------
\subsubsection{Casos de Estudio}
\label{subsubsec:di-casos}
% ---------------------------------------------------------------

Para anclar los estadísticos en viajes concretos, se eligieron doce pares
origen-destino ilustrativos. Ocho de ellos parten de Milpa Alta, la alcaldía con
menor cobertura de la \cdmx{} ($C = 11.75$\,\%) y la más alejada de la red de
transporte estructurado; los demás tienen origen en otros puntos de la periferia
sur y oriente. Los destinos son puntos de referencia conocidos, entre ellos
polos de empleo y educación como Paseo de la Reforma, Santa Fe, Ciudad
Universitaria o la FES Acatlán. Los casos no forman una muestra representativa;
sirven para mostrar trayectos concretos, mientras que la lectura estadística
corresponde a la muestra aleatoria de 100 pares. El Cuadro~\ref{tab:di-casos}
consolida los valores de $DI$ para cada caso.

\begin{table}[H]
    \centering
    \caption[Factor de Desviación — doce casos de estudio]{%
    Factor de Desviación ($DI$) para doce pares origen-destino
    ilustrativos de la \zmvm{}, con origen en las periferias sur y
    oriente; no constituyen una muestra representativa. Corrida: 2026-10-04.}
    \label{tab:di-casos}
    \small
    \begin{tabularx}{\textwidth}{|l|X|X|r|r|r|}
        \hline
        \textbf{\#} & \textbf{Origen} & \textbf{Destino} &
        \textbf{$D_{\text{red}}$ (km)} &
        \textbf{$D_{\text{Hav.}}$ (km)} &
        \textbf{$DI$} \\
        \hline
         1 & Milpa Alta        & Autódromo / Centro-Oriente   & 32.03 & 22.59 & 1.42 \\
         2 & Milpa Alta        & Paseo de la Reforma (CBD)    & 40.94 & 27.50 & 1.49 \\
         3 & Milpa Alta        & Santa Fe (CBD)               & 41.28 & 28.62 & 1.44 \\
         4 & Milpa Alta        & FES Acatlán                  & 52.40 & 35.73 & 1.47 \\
         5 & Milpa Alta        & Ciudad Universitaria         & 29.30 & 19.96 & 1.47 \\
         6 & Milpa Alta        & Coyoacán                     & 28.40 & 18.86 & 1.51 \\
         7 & Santa Catarina    & Ciudad Universitaria         & 43.13 & 23.10 & 1.87 \\
         8 & Santa Catarina    & FES Acatlán                  & 47.34 & 32.16 & 1.47 \\
         9 & Milpa Alta        & Santa Catarina (Sur-Oriente) & 44.30 & 14.79 & 2.99 \\
        10 & Milpa Alta        & Ajusco                       & 33.01 & 20.56 & 1.61 \\
        11 & Ajusco            & Preparatoria 1               & 16.96 & 11.16 & 1.52 \\
        12 & Parque Los Olivos & Centro de Tláhuac            &  2.21 &  2.08 & 1.06 \\
        \hline
        \multicolumn{5}{|l|}{\textbf{Mejor caso de la muestra\textsuperscript{b}:} Plutarco Elías Calles $\to$ Fuente de Loreto} & \textbf{1.03} \\
        \multicolumn{5}{|l|}{\textbf{Peor caso de la muestra:} Av. Sur del Comercio $\to$ C. Violeta y Nezahualcóyotl} & \textbf{4.42} \\
        \hline
    \end{tabularx}
    \vspace{0.4em}

    \noindent\textsuperscript{b}\,Empata en $DI = 1.03$ con el par
    Guadalupe Inclán Ramírez - Periférico $\to$ Periférico - Tec de Monterrey,
    de los Corredores Concesionados.

    \fuentefigura{Elaboración propia con \texttt{VFTModel}
    \citep{cabrera2026VFTModel}. Corrida: 2026-10-04.}
\end{table}

Entre los doce casos, el Caso~9 tiene el $DI$ más alto ($DI = 2.99$). Milpa Alta
y Santa Catarina Yecahuizotl, ambos en el sur-oriente, están a 14.8~km en línea
recta, pero la red obliga a recorrer 44.3~km entre ellos: casi el triple. No es
un accidente de la muestra, sino la medida de la falta de conectividad
transversal en ese cuadrante de la cuenca.

% ---------------------------------------------------------------
\subsubsection{Vista Integral de la Red}
\label{subsubsec:di-vista-integral}
% ---------------------------------------------------------------

¿Cómo se distribuye el costo del rodeo sobre el territorio? La
Figura~\ref{fig:dashboard-di} responde esa pregunta articulando en un solo panel la
magnitud del $DI$ por alcaldía, la dispersión estadística de la muestra y la
huella geográfica de las rutas analizadas.

\begin{figure}[H]
    \centering
    \includegraphics[width=0.95\textwidth]{%
        Figures/Mapas/Tableau_Resources/RedActual/VFTModel/Dashboards/%
        Dashboard_VFTModel_2026_Factor_desviacion}
    \caption[Dashboard del Factor de Desviación — red multimodal
    \zmvm{}]{%
    Vista integral del Factor de Desviación ($DI$) sobre la red de
    transporte público de la \zmvm{}, generada con el
    \termino{dashboard} analítico del \texttt{VFTModel}.
    Corrida: 2026-10-04.}
    \label{fig:dashboard-di}
    \fuentefigura{Elaboración propia con \textbf{Tableau};
    datos: \texttt{VFTModel} \citep{cabrera2026VFTModel}.}
\end{figure}

Entre las alcaldías de la \cdmx{}, Tláhuac registra el $DI$ medio más alto del
Cuadro~\ref{tab:di-alcaldia} (3.26), aunque sobre solo dos rutas, por lo que el
valor es indicativo. El Caso~7 lo ilustra con un trayecto concreto: de Santa
Catarina Yecahuizotl a Ciudad Universitaria hay 23~km en línea recta, pero la
red exige 43~km de recorrido real ($DI = 1.87$), porque ninguna ruta de los
datos procesados por \texttt{Apímetro} conecta esos dos puntos del sur sin
internar el trayecto hacia el centro de la ciudad. No es un caso límite ni un
atípico estadístico: es la geometría de una red sin conectividad transversal en
ese cuadrante, medida con los mismos instrumentos que describen el resto del
capítulo.

\begin{figure}[H]
    \centering
    \includegraphics[width=0.85\textwidth]{%
        Figures/Cap5/FactorDesviacion/Mapa/NB04_mapa10_caso07.png}
    \caption[Caso~7: Santa Catarina Yecahuizotl → Ciudad
    Universitaria ($DI = 1.87$)]{%
    Caso~7: ruta óptima de Santa Catarina Yecahuizotl a Ciudad
    Universitaria. La red obliga a un recorrido de 43.13~km para
    cubrir una distancia lineal de 23.10~km ($DI = 1.87$), al no
    existir conectividad transversal directa en el cuadrante
    sur-oriente de la \zmvm{}. El trazo es ilustrativo y proviene de la
    corrida 2026-04-27; las distancias corresponden a la corrida
    2026-10-04.}
    \label{fig:di-caso07}
    \fuentefigura{Elaboración propia con \texttt{VFTModel}
    \citep{cabrera2026VFTModel}.}
\end{figure}

Los cuatro indicadores de Fase~1 —impedancia, cobertura, Fuerza Capilar y Factor
de Desviación— describen colectivamente una red metropolitana con capacidades
desiguales: bien articulada en el centro, con déficit progresivo hacia la
periferia sur y oriente, y sin conectividad transversal en ese cuadrante. Esa
descripción es topológica: se basa en la estructura del grafo, sin considerar
las condiciones reales de operación. La Fase~2 incorpora esas condiciones
mediante dos parámetros de ponderación dinámica —el Coeficiente de Fricción Vial
($CF$) y la penalización por transferencia ($T_b$)— que ajustan el modelo a la
experiencia de viaje efectiva.
